\section{CLIP：把图像和文本放进同一个空间}

CLIP 是这讲的第一个主角。它的核心是把图像和文本都编码到同一个向量空间里，然后用对比学习把匹配对拉近，把不匹配的样本推远。它不直接学类别标签，而是学跨模态对齐。

\subsection{对比学习的直觉}

CLIP 的训练逻辑，本质上是把“哪个图和哪个文本更像”变成 batch 内的匹配问题。图像编码器和文本编码器各自把输入压成向量，训练时让真实配对的图文相似度更高，其他配对相似度更低。这个目标和 SimCLR 的直觉非常接近，只是从“同一张图的不同增强”扩展成了“图像与语言描述”。\footnote{字幕 00:06:40--00:08:20 中，讲者明确把 CLIP 的图文对比目标放回到 self-supervised contrastive learning 的框架里。}

\begin{importantbox}{CLIP 的训练目标}
训练数据是一大批图文对。模型同时优化两个方向：
\begin{itemize}
    \item 图像向量应最接近自己的文本描述。
    \item 文本向量也应最接近自己的图像。
\end{itemize}
这是一种对称的 batch 内对比损失。
\end{importantbox}

\begin{figure}[H]
\centering
\includegraphics[width=0.95\textwidth]{slides-images/slide-025.jpg}
\caption{图文对比损失的直觉：把配对样本拉近，把其他样本推远}
\end{figure}

\subsection{为什么 scale 会决定 CLIP 的成败}

CLIP 并不是“一个新的小技巧”突然赢了，而是模型规模和数据规模同时上来了。讲者给出两个关键数字：模型规模大约是 307M 参数，训练图文对来自互联网，规模大约到 4 亿对量级。相比 ImageNet 的 1.3M 图像，这已经是完全不同的数据世界。\footnote{字幕 00:14:10--00:16:10 和 00:17:50--00:18:20 对规模因素有直接说明。}

\begin{figure}[H]
\centering
\includegraphics[width=0.95\textwidth]{slides-images/slide-039.jpg}
\caption{CLIP 强的原因之一：模型和数据都大幅扩展}
\end{figure}

\begin{figure}[H]
\centering
\includegraphics[width=0.95\textwidth]{slides-images/slide-040.jpg}
\caption{CLIP 的参数量远大于传统监督分类模型}
\end{figure}

\begin{figure}[H]
\centering
\includegraphics[width=0.95\textwidth]{slides-images/slide-041.jpg}
\caption{CLIP 的训练数据规模也远大于传统监督分类}
\end{figure}

\begin{warningbox}{为什么 scale 不是“可有可无”}
CLIP 的很多能力并不是来自某个神奇结构，而是来自足够大的 batch、足够大的模型和足够多样的图文监督。没有 scale，很多看起来优雅的对比学习结论其实都学不出来。
\end{warningbox}

\subsection{零样本分类怎么做}

CLIP 最有影响力的能力，是 zero-shot 分类。做法很简单：为每个类别写一个文本 prompt，编码成文本向量，再把输入图像编码成图像向量，用最近邻或点积选出最匹配的类别。\footnote{字幕 00:11:00--00:14:10 中，讲者用 plane/dog/bird 的例子解释了 zero-shot 分类。}

\begin{figure}[H]
\centering
\includegraphics[width=0.95\textwidth]{slides-images/slide-026.jpg}
\caption{为每个类别用文本编码器生成一个向量，作为类别原型}
\end{figure}

\begin{figure}[H]
\centering
\includegraphics[width=0.95\textwidth]{slides-images/slide-028.jpg}
\caption{把 CLIP 看成一个 1-NN 分类器：类别原型来自文本向量}
\end{figure}

对新数据集，流程就是：
\begin{enumerate}
    \item 为类别写 prompt，例如 `a photo of a dog`。
    \item 把 prompt 编码成类别原型。
    \item 把图像编码成向量。
    \item 按相似度做分类。
\end{enumerate}

\begin{knowledgebox}{为什么 prompt 很重要}
CLIP 训练时见到的是句子片段，不是孤立词，所以单个类别词往往不如带语境的短语好用。prompt ensemble 进一步把多个描述取平均，通常还能再提升一点性能。
\end{knowledgebox}

\begin{figure}[H]
\centering
\includegraphics[width=0.95\textwidth]{slides-images/slide-030.jpg}
\caption{同一类别用多种 prompt，再取平均向量作为最终类别表示}
\end{figure}

\subsection{为什么 CLIP 会泛化得更好}

讲者对这个问题给了两个解释。第一，互联网文本比 ImageNet label 丰富得多，它不是只有“dog”这样的单词，而是包含形状、颜色、关系、场景和上下文信息。第二，图文对的总量远超过手工分类数据，所以模型见过的组合更广、更杂，也更容易对 out-of-distribution 样本保持鲁棒。\footnote{字幕 00:14:50--00:16:20 和 00:16:20--00:17:40 对 CLIP 的泛化原因做了直接总结。}

\begin{figure}[H]
\centering
\includegraphics[width=0.95\textwidth]{slides-images/slide-045.jpg}
\caption{CLIP-style 模型在 ImageNet 上的表现已经足以作为新一代分类 foundation model}
\end{figure}

\begin{importantbox}{CLIP-style 模型的优点}
\begin{itemize}
    \item 训练目标简单，容易扩展。
    \item 推理很快，检索和分类都方便。
    \item open-vocabulary 能力强，适合跨域迁移。
    \item 很容易和其他模型做 chaining。
\end{itemize}
\end{importantbox}

\subsection{CLIP 的局限：batch、组合关系和分布外样本}

讲者并没有把 CLIP 神化。相反，后半部分花了很多时间讲它的局限。第一个问题是 batch size 依赖：batch 越大，负样本越丰富，模型越有机会学到细粒度差异；但如果 batch 太小，就容易只学到“猫和车不一样”这种粗粒度区分，学不到真正细的类别边界。\footnote{字幕 00:21:50--00:23:40 对 batch size 和 hard negatives 的重要性有明确讨论。}

\begin{figure}[H]
\centering
\includegraphics[width=0.95\textwidth]{slides-images/slide-046.jpg}
\caption{CLIP-style 模型的优势和局限：高效、零样本，但也会被某些语义混淆击中}
\end{figure}

第二个问题是 compositionality。`mug in grass` 和 `grass in mug` 词元相同，但语义完全不同；CLIP 很容易混掉这种关系，因为它更擅长整体对齐，而不是精细的关系解析。

\begin{figure}[H]
\centering
\includegraphics[width=0.95\textwidth]{slides-images/slide-047.jpg}
\caption{CLIP 会混淆语序和组合关系，说明它对 compositionality 仍然不够强}
\end{figure}

\begin{warningbox}{组合关系为什么难}
组合关系要求模型不只识别对象，还要识别对象之间的角色和方向。只靠更大的 batch 能改善一些细粒度区别，但并不能自动学会真正的结构化语义。
\end{warningbox}

\subsection{hard negatives 不是银弹}

一个自然的修补方向是 hard negative fine-tuning：专门把容易混淆的反例放进 batch，逼模型学会更细粒度区别。这个方法短期内确实能改善部分混淆，但讲者也提醒，它可能带来新的问题：过度强调反例会让模型把本来应该接近的样本也强行推开，最后反而损害 generalization。\footnote{字幕 00:23:40--00:24:55 中，讲者明确说 hard negative 训练最终可能伤害泛化。}

\begin{figure}[H]
\centering
\includegraphics[width=0.95\textwidth]{slides-images/slide-049.jpg}
\caption{用 hard negatives 继续微调可以缓解一些混淆问题}
\end{figure}

\begin{figure}[H]
\centering
\includegraphics[width=0.95\textwidth]{slides-images/slide-051.jpg}
\caption{过度依赖 hard negatives 也会制造新的训练偏差}
\end{figure}

\begin{importantbox}{一个更成熟的判断}
CLIP 的问题不是“还不够大”，而是它需要更好的数据构造策略、更明确的组合监督和更稳定的训练信号。简单堆 batch 并不能自动解决语义结构问题。
\end{importantbox}

